\section*{Ex.\ 1.0 --- Three Papers on OS History and Architecture}

\subsection*{Paper 1 --- Corbat\'o and Vyssotsky, ``Introduction and Overview of the Multics System'' (1965)}
\label{p4:sec:multics}

\textbf{Summary.} This paper opens the six-paper Multics session at the 1965 Fall Joint
Computer Conference. Its problem is access: under batch processing, user contact with large
machines had ``retrogressed'', isolating the programmer from cause and effect. The goal is
a \emph{computer utility}, running ``continuously and reliably 7 days a week, 24 hours a
day in a way similar to telephone or power systems''.

The mechanisms are co-designed with the GE~645. Addressing is two-dimensional: programs are
written as segments that grow or shrink during execution, up to a quarter million per user
of a quarter million words each, paged with either 64- or 1{,}024-word pages, and
cross-segment references are bound by dynamic linking at first reference. Pure shareable
procedures are the normal mode, so there is ``no clear-cut demarcation between user
programs and system programs''; protection comes from descriptor bits acting as hardware
``fire-walls''; processors form an anonymous pool in which ``the supervisor does not have a
special processor''; and PL/I is adopted for machine independence. There is no evaluation,
since the system did not yet exist, and the one quantitative claim, ``a few hundred''
simultaneous users by ``simple scaling of processor and memory speed'', is immediately
called ``hazardous''.

\textbf{Critical judgment.} This is a prospectus, and nothing in it is falsifiable by its
own evidence. The deeper omission is not measurement but cost. The paper proposes to meet
``almost all of the present and near-future requirements'' of a large installation and
never asks what that generality would cost to build, which is precisely where the project
failed: development ran years late, Bell Labs withdrew in 1969, and UNIX was written
shortly afterwards in reaction to the scale. Its most accurate sentences are its own
hedges, that the plans are ``not unattainable'' but that it would be ``presumptuous'' to
expect the initial system to meet all of them.

History has split the substance cleanly. Vindicated: paged virtual memory,
device-independent I/O, a symbolically named hierarchical file system with per-file access
control and backup, symmetric multiprocessing with no master, and writing an operating
system in a high-level language. Refuted: user-visible segmentation lost to the flat paged
address space, and x86-64 all but abolished it; dynamic linking survived, but as lazy
binding over a flat space rather than as segment binding.

The computer utility is the interesting case. The economic argument, that ``bulk
economies'' and ``floor space, management efficiency and operating personnel'' favour
``centralizing computer facilities in a single large installation'', was falsified for some
thirty years by the minicomputer and the personal computer, then vindicated by cloud
computing. The paper was right about the destination and wrong about the reason, since what
made centralisation win was statistical multiplexing at datacentre scale.

\subsection*{Paper 2 --- Accetta et al., ``Mach: A New Kernel Foundation for UNIX Development'' (1986)}
\label{p4:sec:mach}

\textbf{Summary.} Mach's target is the accretion of mechanism in UNIX. A system that began
with a small uniform set of operations on file descriptors had grown ``a staggering
number'' of overlapping facilities, from streams and sockets to ``a mind-boggling array of
\texttt{ioctl} operations'', while the process abstraction remained too heavyweight to
express parallelism.

Mach replaces the base with four abstractions. A \emph{task} owns a paged address space and
its port rights; a \emph{thread} is the unit of CPU utilisation, so a UNIX process is a
task with one thread; a \emph{port} is a kernel-protected message queue acting as a
capability, with any number of senders but one receiver; and a \emph{message} is typed data
that may carry port rights and may be as large as an entire address space. Every operation
on any other object is a message to the port representing it, so that ``the actual system
running on any particular machine is a function of its servers rather than its kernel''.

The virtual memory system offers per-page inheritance and a current-versus-maximum
protection pair in which the maximum ``may never be raised, it may only be lowered'', with
copy-on-write implementing both \texttt{fork} and bulk message transfer. Most
consequentially, page faults may be serviced outside the kernel by user-level pagers, so a
memory-mapped file is simply memory whose pager is the file system. As of April 1986 Mach
was binary compatible with 4.3BSD and ran on most VAX machines, a four-processor
VAX~11/784 and the IBM RT/PC.

\textbf{Critical judgment.} The gap between the ambition of the design and the weight of
the evidence is very wide, and the paper is candid that ``extensive performance comparisons
with other systems have not yet been done''. Its entire quantitative evaluation is one
number: on a MicroVAX~II, touching newly allocated memory costs ``less than 0.7
milliseconds per 1024 bytes of data (versus approximately 1.2 milliseconds for 4.3BSD)''.
That microbenchmark measures zero-fill fault cost, exactly what the new VM system was built
to improve, and \texttt{fork} is called ``substantially faster'' with no figure attached at
all.

Nothing measures IPC latency, although ports and messages are the centrepiece of the
design; nothing measures multiprocessor scaling, although the abstract's opening sentence
announces a multiprocessor kernel; threads were not yet implemented and were ``expected by
Summer 1986''; and the caption of Figure~6 concedes that as of April 1986 the box labelled
UNIX compatibility ``still executes in kernel state''. The central claim, that operating
system services can be moved out to user-level servers at acceptable cost, is precisely the
claim the paper does not test.

History then split the verdict. The mechanisms won comprehensively: separating the address
space from the schedulable entity is now universal, memory-mapped files and external pagers
are standard, the machine-dependent split survives as the \texttt{pmap} layer in the BSDs,
and macOS and iOS run on a Mach derivative to this day. The performance thesis lost: Mach
3.0's user-space UNIX server proved slow, Chen and Bershad traced the degradation to
memory-system behaviour rather than message path length, and Liedtke's L4 work argued the
cost was an artefact of implementation rather than of microkernels.

\subsection*{Paper 3 --- Baumann et al., ``The Multikernel: A New OS Architecture for Scalable Multicore Systems'' (2009)}
\label{p4:sec:multikernel}

\textbf{Summary.} The premise is that hardware now diversifies faster than system software
can be retuned for it, so an OS built as a shared-memory kernel with lock-protected data
structures faces an unbounded engineering bill; the authors evidence this with the removal
of the Windows~7 dispatcher lock, which ``touched 6000 lines of code in 58 files'' and was
described as heroic. Arguing that the machine has itself become a network, they propose
three principles: make all inter-core communication explicit, make OS structure
hardware-neutral, and treat state as replicated rather than shared.

Their prototype, Barrelfish, splits the per-core OS instance into a privileged
\emph{CPU driver} (7135 lines of C for x86-64) and a schedulable user-space \emph{monitor}
that holds the replicas and runs agreement protocols over URPC, a message transport tuned
to the cache-coherence protocol. The motivating measurement is stark: a single core updates
shared state in under 30 cycles, but with 16 cores on the same data each update costs
``almost 12{,}000 extra cycles'', and a single dedicated server ``can perform twice as many
updates per second as all 16 shared-memory threads combined''. The evaluation then shows
NUMA-aware multicast shootdown holding near 3.1k cycles across 32 cores where broadcast and
unicast climb to roughly 13.8k and 12.7k; unmapping at 32 cores costing about 17.5k cycles
against 29.5k for Linux and 46k for Windows; and a static web server at 18{,}697 requests
per second against 8924 for lighttpd on Linux.

\textbf{Critical judgment.} This is much the strongest evaluation of the three and an
unusually honest one. It states outright that ``it would be wrong to draw any quantitative
conclusions from our large-scale benchmarks'', concedes that the capability system was ``a
mistake\ldots unnecessarily complex, and no more efficient than scalable perprocessor
memory managers used in conventional OSes like Windows and Linux'', calls its network stack
``very much a placeholder'', and admits every test platform was homogeneous x86-64, so the
heterogeneity motivating the whole argument was never exercised.

That honesty makes the weak points easier to locate. The web-server win pits a
purpose-built user-space lwIP stack against general-purpose lighttpd on Linux with offload
enabled that Barrelfish's driver did not support, so what is measured is specialisation
rather than architecture. The scalability case rests on global-coordination
microbenchmarks, and the gain there comes largely from replacing serial IPIs with a
topology-aware multicast tree, an algorithmic change a shared-memory kernel could equally
adopt, so the experiment does not isolate shared-nothing structure as the cause.

Reading the two figures against each other is more revealing. The raw messaging baseline
costs about 3.1k cycles at 32 cores while the complete unmap costs about 17.5k, so more
than four-fifths of the operation is Barrelfish's own monitor LRPC, marshaling and
unoptimised thread dispatch. The unmap curve makes the same point at the other end:
Barrelfish does not start cheaper, costing about 10.3k cycles at two cores against roughly
5.2k for Linux and 5.8k for Windows, and what it buys is slope, growing about
1.7$\times$ to 32 cores where Linux grows 5.7$\times$ and Windows 7.9$\times$. The
multikernel did not reduce the constant cost of the operation; it roughly doubled it and
bought a much flatter curve. That is the trade the authors concede, that routing
invocations through monitor RPC adds ``a constant overhead on current hardware of several
thousand cycles'' which is ``constant as the number of cores increases'', a bargain that
pays only once the machine is wide enough and whose break-even against Linux does not
arrive until around twelve or thirteen cores. The messaging table reinforces the point,
showing URPC at 450 cycles of latency against 424 for L4's decade-older IPC, the real gains
being throughput and cache footprint rather than latency.


\subsection*{References}
\label{p4:sec:refs}

\begin{enumerate}
  \item[{[1]}] F.~J. Corbat\'o and V.~A. Vyssotsky. ``Introduction and Overview of the
    Multics System.'' In \emph{Proceedings of the AFIPS Fall Joint Computer Conference
    (FJCC)}, vol.~27, part~1, pp.~185--196, 1965.
  \item[{[2]}] M. Accetta, R. Baron, W. Bolosky, D. Golub, R. Rashid, A. Tevanian and
    M. Young. ``Mach: A New Kernel Foundation for UNIX Development.'' In \emph{Proceedings
    of the USENIX Summer Conference}, pp.~93--112, 1986.
  \item[{[3]}] A. Baumann, P. Barham, P.-E. Dagand, T. Harris, R. Isaacs, S. Peter,
    T. Roscoe, A. Sch\"upbach and A. Singhania. ``The Multikernel: A New OS Architecture for
    Scalable Multicore Systems.'' In \emph{Proceedings of the 22nd ACM Symposium on Operating
    Systems Principles (SOSP)}, pp.~29--44, 2009.
\end{enumerate}
